\section{Heavy elements ($Z\gtrsim$30): Modifications for scalar relativity}
\label{Sec:rel-example}

With H and O, the preceding, simple H$_2$O testcase involves only
elements so light that relativistic effects on their total energies
can be safely ignored in production calculations. 

The rule of thumb that relativistic effects do not matter much for
quantities derived from total energies (binding energies, geometries)
holds up to elements number $Z$=20-30 (Ca or Zn), depending on the
reqired accuracy. For example, the DFT-LDA lattice parameter for fcc
Cu is 3.54~{\AA} in the non-relativistic case, but 3.52~{\AA} if
relativistic effects are included. In any case, relativistic effects
should be accounted for in accurate calculations \emph{beyond} these
elements. As described in more detail in Ref. \cite{Blum08}, this
process is handled by FHI-aims at different levels of approximation,
with some overhead resulting compared to the simple non-relativistic
case. 

In a nutshell, the physical impact of relativity for heavy elements
\emph{is} noticeable not just as a total-energy offset, but
importantly in total energy differences, and in particular impacts
also geometries. The underlying reason is that core and 
valence orbitals ``see'' the near-nuclear region (where relativistic
effects are most important) differently, depending on their 
angular momentum $l$. Somewhat simplistically put, the increased
relativistic mass of the electron near the nucleus results in a
tendency for all orbitals to ``shrink'' compared to the
non-relativistic case, but to a different degree for different
orbitals. The shrinkage thus changes not just atom sizes, but also the
nature of bonding itself. A detailed discussion of relativistic
effects is beyond the scope of this manual (and can be found in many
excellent reviews, e.g., Ref. \cite{Pyykko88}) -- but bear in mind
that relativistic effects should not simply be shrugged off!

As outlined in Ref. \cite{Blum08} and Section \ref{Sec:rel}, FHI-aims
provides two levels of approximation to scalar relativistic effects:
the ``atomic'' and ``scaled'' zero-order regular approximation
(ZORA). Both are invoked by the \keyword{relativistic} keyword. In our
experience, energy differences and geometries from \emph{both}
approximations closely match relativistic benchmark calculations
performed using the full-potential linearized augmented plane wave
method. While the absolute total energies from scaled ZORA are
significantly closer to the fully relativistic limit than atomic
ZORA, atomic ZORA appears to be slightly closer to other benchmark
codes for energy differences, and the implementation of total energy 
\emph{gradients} is much simpler in the atomic ZORA case.

We therefore recommend the following strategy to perform standard
relativistic calculations with FHI-aims:
\begin{itemize}
  \item For forces, geometry changes
    (structure optimization, molecular dynamics, vibrations), energy
    differences, and valence eigenvalues, the simple recommended
    option is \\ ``\keyword{relativistic} \texttt{atomic\_zora} 
    \texttt{scalar}''. 
  \item Optionally (for nicer-looking total energies and core
    eigenvalues), perform a follow-up calculation (fixed geometry)
    using \\ ``\keyword{relativistic} \texttt{zora} 
    \texttt{scalar} \option{threshold}'' \\ (where threshold is a small
    number, e.g., 10$^{-12}$).
\end{itemize}
Note that this reflects a change from earlier recommendations in this
manual; for many years, we were uncertain how (if at all) to qualify
the very small differences between atomic ZORA and scaled ZORA
compared to other benchmark approaches. Following the ``Delta
Project'' (https://molmod.ugent.be/deltacodesdft), it appears that
atomic ZORA is slightly more reliable -- but most importantly, it is
simpler. 

If, for any reason, you do choose to follow up with ``scaled ZORA''
anyway, we emphasize two important points for the practical use of the
``scaled ZORA'' approximation:
\begin{enumerate}
  \item \textbf{Only use the \emph{final} total energy from a converged
    scaled-ZORA run}, given after the scaled ZORA post-processing step
    is complete. During the s.c.f. cycle, the Kohn-Sham wave function
    and all intermediate total energies are obtained at the simple
    ZORA level (\emph{not} the same as atomic ZORA) only, so \emph{any
      intermediate total energy is not reliable}. 
  \item \textbf{For energy differences between different structures,
    use the same \subkeyword{species}{radial\_base} integration
    grids}. Relativistic effects near the nucleus need to be
    integrated accurately, but residual integration errors cancel on
    an atom-by-atom basis. In energy differences, the remaining errors
    for the same \subkeyword{species}{radial\_base} and
    \subkeyword{species}{radial\_multiplier} settings thus cancel out.
\end{enumerate}

We illustrate the practical use of of atomic and (if needed) scaled ZORA for the
Au dimer in DFT-LDA, found in the \emph{testcases/Au\_dimer}
directory. 

In the subdirectory \emph{relax\_light}, a quick but
sufficiently accurate prerelaxation is set up, at the \texttt{atomic\_zora}
level. The \texttt{control.in}
file is set up to use \keyword{relativistic} \texttt{atomic\_zora}
\texttt{scalar} and \keyword{relax\_geometry} \texttt{bfgs}
\texttt{1.e-2} to converge forces down to 10$^{-2}$ eV/{\AA}. Since
this is intended to be a quick prerelaxation run (starting with an
arbitrary separation of 3 {\AA} of both atoms), the ``light'' species 
default settings for Au are used for the \keyword{species}
subsettings. Compared to the much more accurate ``tight'' settings,
\begin{itemize}
\item the radial and angular integration grids are significantly reduced, 
\item the Hartree potential expansion is capped at
  \subkeyword{species}{l\_hartree}=4,
\item the cutoff potential onset and width in \subkeyword{species}{cut\_pot}
  are reduced to a minimum that we still consider safe (3.5 {\AA} / 1.5 {\AA},
  respectively),
\item the basis set is the \emph{spdf} section of \emph{tier} 1 only.
\end{itemize}.
Among these settings, the reduction of the basis set to \emph{spdf} has the
biggest impact on the resulting equilibrium geometry, $d$=2.465~{\AA}. Just
adding the $g$ function of \emph{tier} 1 would already reduce the error by
0.013~{\AA}.  

After this relaxation run is complete, a file \texttt{geometry.in.next\_step} is
written out by FHI-aims. This file contains the final, relaxed geometry from
the ``light'' prerelaxation run just performed, as well as the Hessian matrix
estimate created by the \keyword{relax\_geometry} \texttt{bfgs} relaxation algorithm.
This file can serve as an improved guess for \texttt{geometry.in} in the next, 
typically more expensive ``tight'' post-relaxation step. 

In the subdirectory \emph{postrelax\_tier1}, the final geometry from the 
previous step is used as the starting point for an accurate post-relaxation
using the ``tight'' \emph{species\_defaults} settings. Note that,
while all other settings are tightly converged at this level, the basis set
convergence should normally still be tested explicitly. The \emph{tier} 1
level used here for Au is quite sufficient for an accurate result. The binding
distance converges after two additional relaxation steps, at $d$=2.449 {\AA}.

\emph{Optionally:} In the subdirectory \emph{final\_scaled\_zora}, the
fixed post-relaxed 
equilibrium geometry is used to compute the Au$_2$ binding energy at the
scaled ZORA level.  The computation is a single-point run with
the full \emph{tier} 1 basis set, using the rather safe grid and
Hartree potential settings from the ``tight'' species default
file. FHI-aims output files using the same \texttt{control.in} file
are given for both Au$_2$ and a single Au atom, yielding a binding
energy of 3.263~eV for the dimer. Note that this deviates from a possible
\texttt{atomic\_zora} result by less than 0.02~eV, matching our experience
that energy \emph{differences} usually agree very well between
\texttt{atomic\_zora}, scaled \texttt{zora}, and independently with results
other scalar relativistic all-electron results obtained using other
codes. 

In conclusion, this section illustrates how a quick pre-relaxation followed 
by a safe calculation of the relevant energy differences can be
combined. The key point in this endeavour is to ensure that the final 
relaxed geometry is accurate with as little computational overhead as
possible. 
